\section{A Prompt-Independent Word-Level Store}
\label{sec:store}

Table~\ref{tab:store} reports, for the pure-word tasks, accuracy above the nuisance floor at the
word position of the neutral sentence and at the answer position of the neutral and question arms.

\begin{table}[t]
  \centering\footnotesize
  \setlength{\tabcolsep}{4pt}
  \resizebox{\textwidth}{!}{\begin{tabular}{@{}lrrrrrr@{}}
\toprule
task & shots & floor & word, last token & word, span mean & answer, no question & answer, question \\
\midrule
\texttt{first\_letter} & 0 & 0.457 & +0.434 (16) & +0.508 (16) & +0.436 (16) & +0.537 (16) \\
\texttt{first\_letter} & 3 & 0.475 & +0.415 (16) & +0.492 (16) & +0.396 (16) & +0.522 (16) \\
\addlinespace[2pt]
\texttt{last\_letter} & 0 & 0.547 & +0.314 (16) & +0.177 (16) & +0.115 (16) & +0.356 (16) \\
\texttt{last\_letter} & 3 & 0.552 & +0.296 (16) & +0.157 (16) & $-$0.004 (6) & +0.393 (16) \\
\addlinespace[2pt]
\texttt{most\_frequent\_letter} & 0 & 0.442 & +0.135 (16) & +0.165 (16) & +0.084 (14) & +0.283 (16) \\
\texttt{most\_frequent\_letter} & 3 & 0.400 & +0.134 (16) & +0.162 (16) & +0.018 (10) & +0.357 (16) \\
\addlinespace[2pt]
\texttt{word\_length} & 0 & 0.363 & +0.003 (7) & $-$0.048 (2) & +0.029 (14) & +0.193 (12) \\
\texttt{word\_length} & 3 & 0.351 & +0.019 (10) & $-$0.018 (4) & +0.062 (15) & +0.383 (16) \\
\addlinespace[2pt]
\texttt{count\_distinct\_lm} & 0 & 0.216 & +0.077 (16) & +0.088 (16) & +0.072 (16) & +0.079 (16) \\
\texttt{count\_distinct\_lm} & 3 & 0.222 & +0.100 (16) & +0.090 (16) & +0.066 (16) & +0.098 (16) \\
\bottomrule
\end{tabular}}
  \caption{Accuracy above the word-only nuisance floor (held-out layer selection), mean over the 16
    models, with the number of models above the floor in parentheses. ``Word'' columns probe the
    query word's tokens in the neutral sentence (last token, or mean over the word's tokens);
    ``answer'' columns probe the final prompt token. For \texttt{count\_distinct\_lm} the floor
    equals the majority baseline because the letter set, which determines the answer, is excluded
    as an oracle.}
  \label{tab:store}
\end{table}

\paragraph{The store is present in every model.}
The word tokens encode \texttt{first\_letter}, \texttt{last\_letter} and
\texttt{most\_frequent\_letter} above the floor in all 16 models, at both shot counts and under
both reads. The size depends on where the property lives in the word: the last-token read is
stronger for \texttt{last\_letter} ($+0.314$ against $+0.177$ for the span mean at zero shots),
and the span mean is stronger for \texttt{first\_letter} and \texttt{most\_frequent\_letter}.
We take the better read per task as the store in what follows; every conclusion below also holds
under the other read.

\paragraph{It is mostly in the embedding.}
At layer 0 the word tokens already carry 73\% of \texttt{last\_letter}'s above-floor store
(last-token read, zero shots; 69\% at three shots) and 82\% of \texttt{most\_frequent\_letter}'s
(span mean; 75\% at three shots). For \texttt{first\_letter} the share is 53--56\%, with the rest
added by intermediate layers. This agrees with direct evidence that token embeddings encode
spelling \citep{kaushal2022tokens,itzhak2022spelling}. In the present setting it means that the
prompt-independent store is largely lexical: it is a property of the tokens, read off with
little contextual processing.

\paragraph{It does not grow reliably with scale.}
Across the 16 models the store's correlation with $\log_{10}$ parameters is $r = 0.47$
($p = 0.064$) for \texttt{last\_letter}, $r = -0.05$ for \texttt{most\_frequent\_letter} and
$r = -0.16$ for \texttt{first\_letter} (zero shots, last-token read; span-mean values are similar).
By contrast, the question arm's answer position correlates with scale at $r = 0.73$
($p = 0.002$) and $r = 0.72$ ($p = 0.002$) for the first two tasks.

\paragraph{Where there is no store.}
\texttt{word\_length} is not linearly decodable at the word position above its floor, which here
is the letter-set predictor ($+0.003$ last token, $-0.048$ span mean at zero shots). Its answer
position under the question reaches $+0.383$ at three shots, so length, unlike the letter
properties, is assembled at the answer position rather than read off the word.
\texttt{count\_distinct\_lm} has a small store ($+0.08$ to $+0.10$ above majority in all 16 models)
that the question barely changes.

\paragraph{Sensitivity to the preceding context.}
The store is measured after the prefix \texttt{The word "}. In a pre-registered follow-up
(Appendix~\ref{app:prereg}) the prefix \texttt{In the word "} shifted the same word tokens' accuracy
by $+0.02$ to $+0.05$, with the word tokenized identically. Prefix sensitivity of this size is small
next to the store itself ($+0.13$ to $+0.51$ above the floor), but it sets the resolution at which
word-level differences between prompts can be interpreted.
